\section{一页式最终结论}

本章计算目的：先给出最终实验决策答案，使汇报读者在进入计算细节前就能看到 100--1000 m 下各源项强度、排序和实验优先级。

\begin{table}[H]
\centering
\scriptsize
\caption{水下舰艇电场源项在不同距离下的强度范围与实验意义}
\label{tab:final-source-distance-strength}
\resizebox{\textwidth}{!}{
\begin{tabular}{p{2.5cm}p{1.8cm}p{3.2cm}p{3.2cm}p{3.2cm}p{3.2cm}p{1.7cm}p{3.4cm}}
\toprule
源项 & 频率特征 & 100 m & 300 m & 500 m & 1000 m & 是否可远场排序 & 对实验的含义 \\
\midrule
UEP/ICCP 静态或慢变场
& 0 Hz/慢变
& $5.97\times10^{-7}$--$1.19\times10^{-4}$ V/m；$0.597$--$119$ $\mu$V/m
& $2.21\times10^{-8}$--$4.42\times10^{-6}$ V/m；$0.0221$--$4.42$ $\mu$V/m
& $4.77\times10^{-9}$--$9.55\times10^{-7}$ V/m；$0.00477$--$0.955$ $\mu$V/m
& $5.97\times10^{-10}$--$1.19\times10^{-7}$ V/m；$5.97\times10^{-4}$--$0.119$ $\mu$V/m
& 是
& 幅值最大；用于慢变场、空间梯度和多通道一致性验证。 \\
\midrule
轴频基频
& $f_s=3$ Hz
& $5.97\times10^{-9}$--$1.19\times10^{-5}$ V/m；$0.00597$--$11.9$ $\mu$V/m
& $2.21\times10^{-10}$--$4.42\times10^{-7}$ V/m；$2.21\times10^{-4}$--$0.442$ $\mu$V/m
& $4.77\times10^{-11}$--$9.55\times10^{-8}$ V/m；$4.77\times10^{-5}$--$0.0955$ $\mu$V/m
& $5.97\times10^{-12}$--$1.19\times10^{-8}$ V/m；$5.97\times10^{-6}$--$0.0119$ $\mu$V/m
& 是
& 幅值低于静态场，但线谱稳定；当前第一实验目标。 \\
\midrule
轴频高阶谐波
& $2f_s,3f_s=6,9$ Hz
& $1.99\times10^{-9}$--$5.97\times10^{-6}$ V/m；$0.00199$--$5.97$ $\mu$V/m
& $7.37\times10^{-11}$--$2.21\times10^{-7}$ V/m；$7.37\times10^{-5}$--$0.221$ $\mu$V/m
& $1.59\times10^{-11}$--$4.77\times10^{-8}$ V/m；$1.59\times10^{-5}$--$0.0477$ $\mu$V/m
& $1.99\times10^{-12}$--$5.97\times10^{-9}$ V/m；$1.99\times10^{-6}$--$0.00597$ $\mu$V/m
& 是
& 用于验证 6/9 Hz 谱线分辨能力；优先级低于 3 Hz 基频。 \\
\midrule
尾流/流动诱导电场
& 低频宽带
& 局部上限 $\le vB_0$；不可按 100--1000 m 远场排序
& 局部上限 $\le vB_0$；不可按 100--1000 m 远场排序
& 局部上限 $\le vB_0$；不可按 100--1000 m 远场排序
& 局部上限 $\le vB_0$；不可按 100--1000 m 远场排序
& 否
& 后续单独模型；不作为当前主验证目标。 \\
\midrule
50 Hz 工频干扰
& 50 Hz 及倍频
& 由实测 $A_{50,i}/d_{\mathrm{eff}}$ 决定；不可作为真实舰艇源项强度
& 由实测 $A_{50,i}/d_{\mathrm{eff}}$ 决定；不可作为真实舰艇源项强度
& 由实测 $A_{50,i}/d_{\mathrm{eff}}$ 决定；不可作为真实舰艇源项强度
& 由实测 $A_{50,i}/d_{\mathrm{eff}}$ 决定；不可作为真实舰艇源项强度
& 否
& 必须测量和剔除；不解释为舰艇主信号。 \\
\bottomrule
\end{tabular}}
\end{table}

\textbf{幅值最大：} 在 100--1000 m 内，UEP/ICCP 静态或慢变场幅值最大，约比 $m_1=0.03$ 的轴频基频高 $30.5\ \mathrm{dB}$，但缺少稳定线谱特征。

\textbf{最适合验证：} 当前实验最适合先做轴频基频，尤其 100 m 等级链路验证和 300--500 m 噪声底验证，因为 3 Hz 线谱明确。

\textbf{干扰剔除：} 50 Hz 工频只按 $A_{50,i}$ 和谐波实测处理，进入屏蔽、接地、差分和窄带剔除流程，不进入舰艇源项排序。

\textbf{后续问题：} 尾流/流动诱导只保留局部上限，暂不作为主线；若研究尾流，应另建流体-电磁耦合模型。

本章对实验决策的作用：表 \ref{tab:final-source-distance-strength} 是全文唯一主表，后续章节只解释该表如何计算得到。
